%%%%%%%%%%%%%%%%%%%%%%%%%%%%%%%%%%%%%%%%%%%%%%%%%%%%%%%%%%%%%%%%%%%%%%%%%%%
%%  Project: CPHOS LaTeX Template
%%  File: example-problem.tex
%%  Copyright © 2026 gameswu & CPHOS (www.cphos.cn)
%%
%%  Permission is hereby granted, free of charge, to any person obtaining a 
%%  copy of this software and associated documentation files (the "Software"), 
%%  to deal in the Software without restriction, including without limitation 
%%  the rights to use, copy, modify, merge, publish, distribute, sublicense, 
%%  and/or sell copies of the Software, and to permit persons to whom the 
%%  Software is furnished to do so, subject to the following conditions:
%%
%%  The above copyright notice and this permission notice shall be included 
%%  in all copies or substantial portions of the Software.
%%
%%  THE SOFTWARE IS PROVIDED "AS IS", WITHOUT WARRANTY OF ANY KIND, EXPRESS 
%%  OR IMPLIED, INCLUDING BUT NOT LIMITED TO THE WARRANTIES OF MERCHANTABILITY, 
%%  FITNESS FOR A PARTICULAR PURPOSE AND NONINFRINGEMENT. IN NO EVENT SHALL 
%%  THE AUTHORS OR COPYRIGHT HOLDERS BE LIABLE FOR ANY CLAIM, DAMAGES OR OTHER 
%%  LIABILITY, WHETHER IN AN ACTION OF CONTRACT, TORT OR OTHERWISE, ARISING 
%%  FROM, OUT OF OR IN CONNECTION WITH THE SOFTWARE OR THE USE OR OTHER 
%%  DEALINGS IN THE SOFTWARE.
%%%%%%%%%%%%%%%%%%%%%%%%%%%%%%%%%%%%%%%%%%%%%%%%%%%%%%%%%%%%%%%%%%%%%%%%%%%

% CPHOS理论命题模板

\documentclass[answer]{cphos}

% === 设置文档信息 ===
\cphostitle{CPHOS物理竞赛联考}
\cphossubtitle{理论试题命题模板}

\setscorecheck{true}  % 启用分值检查，可设为false关闭

\begin{document}
\begin{problem}[50]{量子比特的制备}

% === 题干 ===
\begin{problemstatement}
量子计算利用量子比特的叠加与纠缠特性，能并行处理处理海量数据，在特定问题上远超经典计算机。自上世纪八十年代该概念产生，到近年“量子霸权”（Quantum Supremacy）的提出，该行业迅猛发展，成为当今物理科研的热门话题。

在超导量子计算中，我们通过微波电路来人工构建二能级系统，即量子比特。本题将引导你推导并理解如何从一个简单的LC谐振腔，一步步演变为目前主流的Transmon量子比特。

在量子化过程中，物理量被视作算符，乘法交换律一般情况下不再成立。可引入算符$\hat{A}$、$\hat{B}$的对易子$\left[ \hat{A},\hat{B}\right]=\hat{A}\hat{B}-\hat{B}\hat{A}$。

\subq{1} LC电路的量子化

考虑一个由理想电容C和理想电感L组成的无耗散LC振荡电路。设电容器上的电荷为$Q$，穿过电感的磁通为$\varPhi$。

\subsubq{1.1}请写出该经典LC电路的哈密顿量$H$（能量表达式），用$\varPhi$、$Q$表示。

\subsubq{1.2}为了对该电路进行量子化，引入共轭变量$Q$和$\varPhi$的对易子$\left[ \hat{\varPhi},\hat{Q}\right]=\mathrm{i}\hslash$。又该电路的哈密顿量与谐振子的哈密顿量形式相似，故引入产生湮灭算符$\hat{a}^\dagger$、$\hat{a}$满足$\left[ \hat{a},\hat{a}^\dagger \right]=1$且$\hat{\varPhi}=\varPhi_0 \left(\hat{a}+\hat{a}^\dagger\right)$与$\hat{Q}=\mathrm{i}Q_0 \left(\hat{a}^\dagger-\hat{a}\right)$，不妨设$\frac{\varPhi_0}{Q_0}=\sqrt{\frac{L}{C}}$,请问$\hat{H}$的表达式，用$\hat{a}$、$\hat{a}^\dagger$及其他已知量表示。

\subsubq{1.3}$\hat{H}$的本征态为该电路的能级，已知$\hat{a}^\dagger \hat{a}$的本征值为$n$（$n=0,1,2\dots$），请直接写出该电路能级$E_n$。

\subsubq{1.4}为与经典计算机保持一致，量子计算机通常构造仅在基态与第一激发态间跃迁的二能级系统作为量子比特。请问上述LC量子电路适合作为量子计算机的基本电路吗？若适合，请给出理由。若不适合，说明应该如何安排能级分布。

\subq{2} 约瑟夫森结（Josephson Junction）的引入

为制作量子计算机，物理学家们在电路中引入约瑟夫森结，约瑟夫森结由两块超导体和中间的极薄绝缘层组成，其满足两个约瑟夫森关系
\begin{equation}
    I=I_c\sin\phi,\quad
    V=\frac{\hslash}{2e} \frac{\mathrm{d}\phi}{\mathrm{d}t}
\end{equation}
其中$\phi$是一个随时间变化的参量。现考虑一个电容与约瑟夫森结形成闭合回路。

\subsubq{2.1} 请计算该电路的哈密顿量（用库伯对数$n_{2e}=\frac{Q}{2e}$与$\phi$及其他已知量表示，常数部分可略去不写）。

\subsubq{2.2} 将该电路的哈密顿量展开至$\phi^4$项，并由此定性说明，为什么约瑟夫森结可作为量子比特被引入。

\subq{3}Transmon量子比特

为解决早期超导量子比特受到环境中电荷噪声$n_{2e}$的严重干扰，物理学家提出了Transmon量子比特，其哈密顿量为
\begin{equation}
    H=4E_C\left(n_{2e}-n_g\right)^2-E_J\cos\phi
\end{equation}
其中$n_g$的引入是为了抑制$n_{2e}$的噪声。

\subsubq{3.1}引入$n_g$在抑制噪声的同时，也带来其自身的噪音。因此，我们该如何设定$E_J, E_C$来有效抑制噪声？

\subsubq{3.2}假设系统参数使得$E_J \gg E_C$，其中$E_C = \frac{e^2}{2C}$ 是充电能。在该极限下，系统的能级间隔对$\phi$的涨落不敏感（transmon regime）。已知在这种情况下，能级可以近似为
\begin{equation}
E_m\approx-E_J+\sqrt{8E_JE_C}\left(m+\frac{1}{2}\right)-\frac{E_C}{12}\left(6m^2+6m+3\right)+\cdots 
\end{equation}
其中$ m = 0,1,2\dots$。计算非谐性（anharmonicity）$\alpha=\left(E_2-E_1\right)-\left(E_1-E_0\right)$，为了使近似更精确，$\frac{E_J}{E_C}$是否可以无限大？

\subq{4}Rabi振荡

\subsubq{4.1}实现量子比特的反转不仅需要能级量子化的系统，还需要外界的驱动。已知在频率$\omega>0$，振幅$a\ll E_1, E_2$的驱动下，由能级$E_1$和$E_2$（$E_1<E_2$）构成的量子比特具有哈密顿算符：
\begin{equation}
	\hat{H}=\begin{pmatrix}
		E_1 & a\cos(\omega t) \\
		a\cos(\omega t) & E_2 
	\end{pmatrix}
\end{equation}
而运动方程写作
\begin{equation}
	\mathrm{i}\hslash\frac{\mathrm{d}}{\mathrm{d}t}\begin{pmatrix}c_1\mathrm{e}^{-\frac{\mathrm{i}E_1t}{\hslash}}\\c_2\mathrm{e}^{-\frac{\mathrm{i}E_2t}{\hslash}}\end{pmatrix}=\hat{H}\begin{pmatrix}c_1\mathrm{e}^{-\frac{\mathrm{i}E_1t}{\hslash}}\\c_2\mathrm{e}^{-\frac{\mathrm{i}E_2t}{\hslash}}\end{pmatrix}
\end{equation}
求出$\frac{\mathrm{d}c_1}{\mathrm{d}t}$，$\frac{\mathrm{d}c_2}{\mathrm{d}t}$。

\subsubq{4.2}由上问结果可以看出，$c_1, c_2$的变化很缓慢，在求平均时将它们看作不变，求出使$\frac{\mathrm{d}c_1}{\mathrm{d}t}$，$\frac{\mathrm{d}c_2}{\mathrm{d}t}$的时间平均值不为$0$的频率$\omega$，此时系统的状态称为Rabi振荡。

\subsubq{4.3}在上述共振频率下，求出量子比特从$|c_1|=1, |c_2|=0$第一次演化到$|c_1|=0, |c_2|=1$所需时间。

\end{problemstatement}

% === 解答 ===
\begin{solution}

\solsubq{1}{23}
\solsubsubq{1.1}{2}
\begin{equation}
   H=\frac{Q^2}{2C}+\frac{\varPhi^2}{2L} \eqtagscore{1}{2}
\end{equation}

\solsubsubq{1.2}{14}
反解出
\begin{equation}
   \hat{a}^\dagger=\frac{1}{2}\left(\frac{\hat{\varPhi}}{\varPhi_0}-\mathrm{i}\frac{\hat{Q}}{Q_0}\right) \eqtagscore{2}{2}
\end{equation}
\begin{equation}
   \hat{a}=\frac{1}{2}\left(\frac{\hat{\varPhi}}{\varPhi_0}+\mathrm{i}\frac{\hat{Q}}{Q_0}\right) \eqtagscore{3}{2}
\end{equation}
代入产生湮灭算符的对易子，可得
\begin{equation}
    Q_0\varPhi_0=\frac{\hslash}{2} \eqtagscore{4}{4}
\end{equation}
结合以
\begin{equation}
    \frac{\varPhi_0}{Q_0}=\sqrt{\frac{L}{C}} \eqtag{5}
\end{equation}
可以解得
\begin{equation}
    \varPhi_0=\sqrt{\frac{\hslash}{2}\sqrt{\frac{L}{C}}},\quad
    Q_0=\sqrt{\frac{\hslash}{2}\sqrt{\frac{C}{L}}} \eqtagscore{6}{2}
\end{equation}
将其代入$\hat{H}$，可以得到
\begin{equation}
    \hat{H}=\hslash\sqrt{\frac{1}{LC}}\left(\hat{a}^\dagger\hat{a}+\frac{1}{2}\right) \eqtagscore{7}{4}
\end{equation}

\solsubsubq{1.3}{2}
\begin{equation}
    E_n=\hslash\sqrt{\frac{1}{LC}}\left(n+\frac{1}{2}\right) \eqtagscore{8}{2}
\end{equation}

\solsubsubq{1.4}{5}
不适合。\addtext{指出不适合}{2}

原因：当我们选择最低的两个能级作为运算空间的时候，其相邻能级之间的能量之差应该不同，否则在外界能量输入时，无法保证原子仅在选取的两个能级间跃迁，运算空间将一片混乱。而LC量子电路体系各能级之差一样，故不适合作为量子计算机的基本电路。\addtext{说明相邻能级差应该不同及原因}{3}

\solsubq{2}{8}
\solsubsubq{2.1}{4}
根据静电能的定义，可得
\begin{equation}
    U=\int_{0}^{t} IV\,\mathrm{d}t=\frac{\varPhi_0I_c}{2\uppi}\left(\cos\phi_0-\cos\phi\right) \eqtagscore{9}{3}
\end{equation}
其中$\varPhi=\frac{h}{2e}$，略去常数，有
\begin{equation}
    H=\frac{2e^2}{C}n_{2e}^2-\frac{\varPhi_0I_c}{2\uppi}\cos\phi \eqtagscore{10}{1}
\end{equation}

\solsubsubq{2.2}{4}
\begin{equation}
    H=\frac{2e^2}{C}n_{2e}^2-\frac{\varPhi_0I_c}{2\uppi}+\frac{1}{2}\frac{\varPhi_0I_c}{2\uppi}\phi^2-\frac{1}{24}\frac{\varPhi_0I_c}{2\uppi}\phi^4 \eqtagscore{11}{2}
\end{equation}
相比于之前的体系，此体系多了恒负的一项。除此之外，我们还可以知道对越高的能级，此项作用越大，故而可知能级差随着能级上升而减小，因此适合作为量子比特。\addtext{解释适合作为量子比特的原因}{2}

\solsubq{3}{6}
\solsubsubq{3.1}{2}
为使噪音足够小，应该让
\begin{equation}
	E_J\gg E_C \eqtagscore{13}{2}
\end{equation}

\solsubsubq{3.2}{4}
代入数值有
\begin{equation}
    \alpha=-E_C\eqtagscore{12}{2}
\end{equation}
能级间隔单调递减，故同上，此系统可以作为量子比特。

$\frac{E_J}{E_C}$不能无限大。\addtext{指出$\frac{E_J}{E_C}$不能无限大}{1}
如果$\frac{E_J}{E_C}$过大，会使得非谐性相对能级差
过小，因此也不可接受。\addtext{说明原因}{1}

\solsubq{4}{13}
\solsubsubq{4.1}{4}
\begin{equation}
	\frac{\mathrm{d}c_1}{\mathrm{d}t}=-\frac{\mathrm{i}ac_2}{\hslash}\cos(\omega t)\exp\left(\frac{\mathrm{i}(E_1-E_2)t}{\hslash}\right)
	\eqtagscore{14}{2}
\end{equation}
\begin{equation}
	\frac{\mathrm{d}c_2}{\mathrm{d}t}=-\frac{\mathrm{i}ac_1}{\hslash}\cos(\omega t)\exp\left(-\frac{\mathrm{i}(E_1-E_2)t}{\hslash}\right)
	\eqtagscore{15}{2}
\end{equation}
\solsubsubq{4.2}{3}
展开
\begin{equation}
	\cos(\omega t)=\frac{\mathrm{e}^{\mathrm{i}\omega t}+\mathrm{e}^{-\mathrm{i}\omega t}}{2} \eqtagscore{16}{1}
\end{equation}
由于当$\varOmega\neq0$时，对时间的平均值
\begin{equation}
	\left\langle\mathrm{e}^{\mathrm{i}\varOmega t}\right\rangle=0
	\eqtag{17}
\end{equation}
所以仅当
\begin{equation}
	\omega=\frac{E_2-E_1}{\hslash} 	\eqtagscore{18}{2}
\end{equation}
时，有非零的平均值。

\solsubsubq{4.3}{6}
\begin{equation}
	\left\langle\frac{\mathrm{d}c_1}{\mathrm{d}t}\right\rangle=-\frac{\mathrm{i}ac_2}{2\hslash}
	\eqtagscore{19}{1}
\end{equation}
\begin{equation}
	\left\langle\frac{\mathrm{d}c_2}{\mathrm{d}t}\right\rangle=-\frac{\mathrm{i}ac_1}{2\hslash}
	\eqtagscore{20}{1}
\end{equation}
因此$c_1, c_2$振荡的频率为
\begin{equation}
	\varOmega=\frac{a}{2\hslash} \eqtagscore{21}{2}
\end{equation}
初始$|c_1|=1, |c_2|=0$，第一次演化到$|c_1|=0, |c_2|=1$时相位为$\frac{\uppi}{2}$，则经过的时间：
\begin{equation}
	t=\frac{\uppi}{2\varOmega}=\frac{\uppi\hslash}{a} \eqtagscore{22}{2}
\end{equation}


\scoring  % 评分标准自动使用题目环境中声明的总分
\end{solution}

\end{problem}

\end{document}